\section{Phương pháp}\label{sec:method}

\subsection{Thiết kế nghiên cứu}

Nghiên cứu có hai phần (Hình~\ref{fig:design}). Phần hồi cứu đọc lịch sử git của ViPharmind để trả lời RQ1
và RQ2. Phần thực nghiệm lấy nhiệm vụ từ commit thật của backend, giấu các test mà commit gốc thêm vào, rồi
giao cùng một đề bài cho tác tử theo hai cách làm chỉ khác nhau ở lời dặn về quy trình. Mỗi tổ hợp (nhiệm vụ,
cách làm, nhóm chạy) chạy ba lần và được chấm tự động bằng test ẩn.

\begin{figure}[H]
\centering
\resizebox{\linewidth}{!}{%
\begin{tikzpicture}[font=\footnotesize, node distance=4mm and 5mm,
  box/.style={draw, rounded corners=2pt, align=center, minimum height=9mm, text width=24mm, fill=gray!6},
  arr/.style={-latex, thick}]
\node[box, text width=22mm] (repo) {Kho mã\\ViPharmind};
\node[box, right=of repo, yshift=9mm] (retro) {Phân tích hồi cứu\\(RQ1, RQ2)};
\node[box, right=of repo, yshift=-9mm] (tasks) {Nhiệm vụ ba tầng\\đề bài + test ẩn};
\node[box, right=of tasks, text width=27mm] (arms) {Hai cách làm\\đặc tả / vibe\\$\times$ nhóm $\times$ ba lần};
\node[box, right=of arms] (grade) {Bộ chấm\\test ẩn, token,\\lỗi tĩnh};
\node[box, right=of grade, text width=20mm] (ana) {Phân tích\\đã đăng ký\\(RQ3--RQ6)};
\draw[arr] (repo.east) -- ++(2mm,0) |- (retro.west);
\draw[arr] (repo.east) -- ++(2mm,0) |- (tasks.west);
\draw[arr] (tasks) -- (arms);
\draw[arr] (arms) -- (grade);
\draw[arr] (grade) -- (ana);
\end{tikzpicture}}
\caption{Thiết kế nghiên cứu: phần hồi cứu và phần thực nghiệm dùng chung một kho mã}\label{fig:design}
\end{figure}

Năm giả thuyết được ghi nguyên văn trong bản đăng ký trước (mục~\ref{sec:prereg}):
\begin{itemize}
\item \textbf{H1} (RQ3). Làm theo đặc tả có tỉ lệ đạt nhiệm vụ cao hơn vibe coding.
\item \textbf{H2} (RQ3). Làm theo đặc tả tốn nhiều token và thời gian hơn.
\item \textbf{H3} (RQ3). Làm theo đặc tả ít lỗi tĩnh, ít hồi quy hơn và có chỉ số bảo trì không kém hơn.
\item \textbf{H4} (RQ5). Khoảng cách giữa hai cách làm ở Sonnet lớn hơn ở Opus, tức quy trình giúp mô hình yếu
hơn nhiều hơn.
\item \textbf{H5} (RQ6). Hiệu số tỉ lệ đạt (đặc tả trừ vibe) tăng theo kích thước nhiệm vụ.
\end{itemize}
Phép thử xác nhận duy nhất là H1 trên nhóm Claude Sonnet ở tầng nhỏ. H5 là phép thử phụ quan trọng nhất, vì
nó kiểm lập luận rằng chi phí cố định của việc viết thiết kế chỉ đáng khi nhiệm vụ đủ lớn. Các phép so còn
lại là phụ.

Trong phần hồi cứu, kho mã ViPharmind được đọc ở hai mốc cố định: nhánh \texttt{main} tại \texttt{3baae61d764f} (sản phẩm) và
nhánh \texttt{camera-ready} tại \texttt{61d10db85f0f} (nghiên cứu). Tập commit là hợp của hai mốc, gồm
\NumCommits{} commit (\NumCommitsFeat{} \texttt{feat}, \NumCommitsFix{} \texttt{fix},
\NumCommitsRefactor{} \texttt{refactor}). Mỗi commit được gán vào khâu có nhiều dòng thay đổi nhất theo luật
đường dẫn file. \NumSpecPlanFiles{} file spec và plan được lấy lại từ lịch sử git. Một commit \texttt{feat}
được coi là \emph{có plan đi trước} khi nó chạm một file mã được nêu trong một plan thêm vào trước nó không quá
14 ngày; quy tắc này được kiểm tay trên \NumLinkChecked{} commit ngẫu nhiên. Mức sửa lại (RQ2) là số commit
\texttt{fix} hoặc \texttt{refactor} chạm cùng file trong 14 ngày sau mỗi commit \texttt{feat}.

\subsection{Nhiệm vụ và hai cách làm}\label{sec:tasks}

Mỗi nhiệm vụ gồm mã nền, đề bài và test ẩn. Ở \emph{tầng nhỏ}, nhiệm vụ là một commit: mã nền là commit cha,
test ẩn là file test mà commit thêm hoặc sửa. Ở \emph{tầng vừa} và \emph{tầng lớn}, nhiệm vụ là một đoạn
commit liền nhau của cùng một plan, test ẩn là mọi file \texttt{test\_*.py} thêm hoặc sửa trong đoạn. Tác tử
không thấy file test mới và chỉ thấy bản cũ của file bị sửa.

Ứng viên tầng nhỏ là mọi commit \texttt{feat} chạm cả \texttt{backend/src} và \texttt{backend/tests}. Bộ chạy
áp tự động các tiêu chí sàng lọc, trong đó tiêu chí ``test ẩn hỏng ở mã nền, đạt ở commit gốc'' được kiểm bằng
cách chạy test thật ở hai commit (Phụ lục~\ref{app:tasks}). Tiêu chí còn lại đòi test ẩn kiểm hành vi qua
giao diện mà đề bài nêu được, nên cần người đọc mã. Thứ tự duyệt \NumEligibleStratifiable{} commit qua sàng
lọc được xáo ngẫu nhiên (hạt giống 20261115) với hạn mức theo lớp mã (\texttt{api} \NumTasksApi{},
\texttt{application} \NumTasksApplication{}, \texttt{domain} \NumTasksDomain{}, \texttt{infrastructure}
\NumTasksInfrastructure{}). Phần thay đổi ở \texttt{backend/src} dài từ \NumTaskSrcLinesMin{} đến
\NumTaskSrcLinesMax{} dòng. Chạy thử cho thấy Claude gần đạt trần ở tầng này với cả hai cách làm, nên nghiên
cứu thêm hai tầng lớn hơn, cắt từ 16 đoạn commit của các plan backend: tầng vừa \texttt{m01}--\texttt{m05}
(\NumMediumSrcLinesMin{} đến \NumMediumSrcLinesMax{} dòng) và tầng lớn \texttt{f08}, \texttt{f14},
\texttt{f15} (\NumFeatureSrcLinesMin{} đến \NumFeatureSrcLinesMax{} dòng).

Đề bài viết bằng tiếng Anh, từ năm đến mười câu, nêu việc cần làm và giao diện phải giữ, không nêu cách cài
đặt và không nhắc tới test. Trợ lý AI soạn nháp, tác giả dự án duyệt. Một lượt đối chiếu tự động mọi tên công
khai mà test ẩn dùng với mã nền và đề bài tìm ra tên thiếu ở ba đề (\texttt{t05}, \texttt{t07},
\texttt{t13}); ba đề được bổ sung trước khi chốt. Đề tầng vừa và lớn có thêm mục ``giao diện bắt buộc'' và
được một tác tử AI độc lập phản biện trước khi chốt. Danh sách đủ nhiệm vụ ở Phụ lục~\ref{app:tasks}.

Hai cách làm nhận cùng đề bài, mô hình, công cụ và giới hạn thời gian; chỉ khác lời dặn đặt trước đề bài
(Phụ lục~\ref{app:prompts}). Lời dặn \emph{vibe coding} chỉ có một câu: cài đặt thay đổi sau đây. Lời dặn
\emph{theo đặc tả} yêu cầu, trong cùng một phiên, viết \texttt{docs/design.md} (mục tiêu, phạm vi, giao diện,
cách làm, cách kiểm thử), viết \texttt{docs/plan.md} chia việc thành các phần, rồi cài đặt từng phần theo kiểu
viết test trước và chạy cả bộ test trước khi dừng. Lượt theo đặc tả không để lại cả hai file được ghi là không
tuân thủ nhưng vẫn được tính (phân tích theo phân bổ, \emph{intention-to-treat}).

Lời dặn theo đặc tả đi theo quy trình Superpowers mà tác giả dự án dùng khi làm ViPharmind, nhưng viết lại
thành một khối chữ thay vì nạp skill (Bảng~\ref{tab:superpowers}). Lý do là các lượt chạy không có người trả
lời, nên bước hỏi đáp không làm được, và cùng một quy trình phải chạy y hệt trên hai công cụ khác nhau. Can
thiệp vì vậy là một gói ``quy trình đặc tả tự động'', không có bước người duyệt.

\begin{table}[H]
\centering\footnotesize
\caption{Lời dặn theo đặc tả so với quy trình Superpowers}
\label{tab:superpowers}
\begin{tabular}{>{\raggedright\arraybackslash}p{0.36\linewidth}>{\raggedright\arraybackslash}p{0.25\linewidth}>{\raggedright\arraybackslash}p{0.29\linewidth}}
\toprule
Bước của Superpowers & Trong lời dặn & Khác biệt \\
\midrule
\texttt{brainstorming} & viết \texttt{docs/design.md} & giữ bản thiết kế, bỏ phần hỏi đáp với người dùng \\
\texttt{writing-plans} & viết \texttt{docs/plan.md} & cùng cấu trúc; Superpowers chia task 2--5 phút kèm mã, lời
dặn để tác tử tự định cỡ task \\
\texttt{test-driven-development} & viết test, thấy test trượt, rồi viết mã & giống \\
\texttt{verification-before-completion} & chạy cả bộ \texttt{uv run pytest} trước khi dừng & giống \\
\texttt{executing-plans},\newline \texttt{subagent-driven-development} & một tác tử làm lần lượt từng task & không
có tác tử con, không có bước duyệt giữa các task \\
\bottomrule
\end{tabular}
\end{table}

\subsection{Điều kiện chạy và chấm}\label{sec:conditions}

Bản đăng ký có bốn nhóm chạy (Bảng~\ref{tab:cohorts}), đều ở mức suy luận medium. Ba lần lặp, hai cách làm và
bốn nhóm cho \NumRunsSmallTotal{} lượt ở tầng nhỏ, \NumRunsMediumTotal{} ở tầng vừa và \NumRunsLargeTotal{} ở
tầng lớn, tổng \NumRunsTotal{} lượt. Nhóm Antigravity hết hạn mức trước khi chạy xong và bị loại
(mục~\ref{sec:data-collected}).

\begin{table}[H]
\centering\small
\caption{Các nhóm chạy}\label{tab:cohorts}
\begin{tabularx}{\linewidth}{@{}lllX@{}}
\toprule
Nhóm & Công cụ (phiên bản) & Mô hình & Vai trò \\
\midrule
\texttt{claude-sonnet} & Claude Code 2.1.325 & Claude Sonnet 5.5 & H1, RQ3, RQ5 \\
\texttt{claude-opus} & Claude Code 2.1.325 & Claude Opus 5.5 & RQ3, RQ5 \\
\texttt{codex-luna} & Codex CLI 0.162.1 & \texttt{gpt-6-luna} & RQ4 \\
\texttt{agy-gemini} & Antigravity CLI 1.5.9 & Gemini 3.8 Flash Medium & RQ4 (bị loại) \\
\bottomrule
\end{tabularx}
\end{table}

Mỗi công cụ chạy agent mặc định với bộ công cụ mặc định, kể cả công cụ web, ở chế độ không tương tác. Mỗi
lượt có một thư mục \texttt{HOME} mới chỉ chứa thông tin đăng nhập, không skill, plugin, bộ nhớ hay file hướng
dẫn riêng (Phụ lục~\ref{app:conditions}). Giới hạn thời gian là 3600 giây ở tầng nhỏ, gấp đôi ở tầng vừa, gấp
ba ở tầng lớn. Thứ tự lượt được xáo ngẫu nhiên để hai cách làm xen kẽ theo thời gian. Lượt bị ngắt vì hạn mức
được chạy lại từ đầu; quyết định chạy lại chỉ dựa trên trạng thái, không dựa trên kết quả.

Để chống rò rỉ lời giải, cây mã giao cho tác tử không có lịch sử git và đã bỏ mọi spec, plan cũ; bộ chạy
từ chối chạy nếu cây còn sót. Mạng mở là hạn chế được chấp nhận: bộ chạy ghi mọi đường dẫn tới repo gốc xuất
hiện trong nhật ký (\texttt{upstream\_urls}) để phân tích độ nhạy.

Trên một bản sao đã bỏ hai file đặc tả, bộ chấm đưa mọi file test có sẵn về bản của mã nền,
chép test ẩn vào và chạy toàn bộ test không mang nhãn \texttt{integration}; tác tử vì vậy không thể ``đạt''
bằng cách sửa test. Bộ chấm cũng đo lỗi tĩnh (\texttt{ruff}, \texttt{pyrefly}), chỉ số bảo trì
(\texttt{radon}), test do tác tử viết và token.

\subsection{Phân tích}\label{sec:analysis}

Chỉ số chính là \emph{đạt nhiệm vụ}: mọi test ẩn đạt và không test cũ nào hỏng. Chỉ số phụ gồm tỉ lệ test ẩn
đạt, hồi quy, lỗi tĩnh, chỉ số bảo trì, test do tác tử viết, token, thời gian và chi phí quy đổi theo bảng giá
API công khai chốt ngày \mbox{2026-11-15}.

Đơn vị phân tích là nhiệm vụ. Hiệu số của một nhiệm vụ là tỉ lệ đạt theo đặc tả trừ tỉ lệ đạt vibe trên các lần
lặp, tính bằng điểm phần trăm (đpt); ước lượng của một nhóm là trung bình hiệu số trên các nhiệm vụ. Khoảng tin
cậy 95\% (KTC) là bootstrap lấy mẫu lại nhiệm vụ (\NumResBootstrapB{} lần). Một giả thuyết có hướng được đọc là
\emph{được ủng hộ} khi KTC nằm hẳn về phía giả thuyết, \emph{trái} khi nằm hẳn về phía ngược lại, còn lại là
\emph{chưa kết luận được}. Nghiên cứu không dùng ngưỡng $p$ và không hiệu chỉnh so sánh nhiều lần, nhưng báo
cáo số phép so.

H4 dùng thiết kế $2\times 2$ (cách làm $\times$ mô hình) và hiệu của hai hiệu số. H5 dùng hiệu của hai hiệu số
(tầng lớn trừ tầng nhỏ), gộp các nhóm, bổ trợ bằng hồi quy logistic hỗn hợp
\texttt{pass \textasciitilde{} arm} $\times$ \texttt{log(src\_lines) + cohort + (1 + arm | task)} với KTC
bootstrap cụm (\NumResModelBootstrapB{} lần). H5 chỉ được ủng hộ khi cả hai cùng hướng, để không kết luận chỉ vì
tầng nhỏ chạm trần. Phân tích độ nhạy chạy lại H1 và H5 sau khi bỏ lượt có \texttt{upstream\_urls}, lượt sửa
cấu hình \texttt{pytest}, hoặc cặp nhiệm vụ chung commit.

Để tìm nguyên nhân thất bại, hai thành viên (Lê Xuân Lợi, Trần Ngọc Bảo) gán nhãn độc lập mọi lượt
không đạt của hai nhóm Claude theo chín nhãn, trên bản vá đã bỏ file đặc tả và mang mã ngẫu nhiên; tác giả dự
án phân xử chỗ lệch. Độ đồng thuận đo bằng kappa của Cohen. Hai người cũng đọc mù \NumPassCheckRuns{} lượt đạt
ngẫu nhiên để tìm lượt đạt test nhưng sai. Ca nghiên cứu định tính là nhiệm vụ tầng lớn có hiệu số lớn nhất ở
\texttt{claude-sonnet}.

\label{sec:prereg}Thiết kế được chốt bằng ba thẻ git: \texttt{prereg-v1} (giả thuyết, chỉ số, kế hoạch phân tích, tầng nhỏ),
\texttt{prereg-v1.1} (nhiệm vụ và đề bài tầng vừa, tầng lớn, trước mọi lượt AI trên hai tầng này) và
\texttt{prereg-v1-analysis} (gói phân tích, đã chạy được trên dữ liệu chạy thử). Không kết quả chính nào được
mở trước thẻ phân tích. Mọi thay đổi sau khi chốt được ghi kèm lý do và những gì đã xem
(Phụ lục~\ref{app:deviations}). Một tác tử AI chạy tách khỏi phiên soạn thảo (Fable 5.2) phản biện bản đăng
ký, đề bài và mã phân tích khi chưa có kết quả.

Bản bổ sung \texttt{prereg-v1.2}, viết \emph{sau} khi mở kết quả chính, thêm hai thí nghiệm thăm dò. E2 cho
mỗi nhiệm vụ thêm một đề sơ sài một hai câu, không nêu tên hàm hay chữ ký, tạo thiết kế $2\times 2$ (mức chi
tiết của đề $\times$ cách làm). Vì test ẩn gọi đúng những tên mà đề sơ sài không nêu, E2 chấm bằng \emph{điểm
danh sách hành vi}: tỉ lệ mục hành vi rút từ test ẩn mà bài đạt, do Fable 5.2 chấm khi không biết cách làm, và
được kiểm bằng chấm tay của hai người. E1 lặp thiết kế gốc với Claude Haiku 5.5. Giả thuyết thăm dò, nêu sau
khi thấy kết quả chính: giá trị của quy trình tăng khi đề càng ít chi tiết. Kết quả E1, E2 chỉ đọc theo dấu và
độ lớn, không dùng để xác nhận H1--H5.
